% Figure from MATH 590 notes, section 11. Compile with the notes preamble.
\begin{tikzpicture}
  \fill[figB!12] (0,0) circle (1.2);
  \draw[figB, thick, dashed] (0,0) circle (1.2);
  \fill[figR!15] (2.4,0) circle (1.2);
  \draw[figR, thick, dashed] (2.4,0) circle (1.2);
  \draw[black!55, thin] (0,0) -- (2.4,0);
  \fill[black] (0,0) circle (1.4pt) node[below=1pt, font=\scriptsize] {$x$};
  \fill[black] (2.4,0) circle (1.4pt) node[below=1pt, font=\scriptsize] {$y$};
  \draw[black!70, fill=white] (1.2,0) circle (1.5pt);
  \draw[decorate, decoration={brace, amplitude=4pt}, black!55] (0,0.12) -- (1.14,0.12) node[midway, above=4pt, font=\scriptsize, black!70] {$\varepsilon$};
  \draw[decorate, decoration={brace, amplitude=4pt}, black!55] (1.26,0.12) -- (2.4,0.12) node[midway, above=4pt, font=\scriptsize, black!70] {$\varepsilon$};
  \draw[decorate, decoration={brace, amplitude=5pt, mirror}, black!55] (0,-1.35) -- (2.4,-1.35) node[midway, below=5pt, font=\scriptsize, black!70] {$d(x,y) = 2\varepsilon$};
  \node[font=\scriptsize, figB, anchor=south east] at (-0.88,0.88) {$B(x,\varepsilon)$};
  \node[font=\scriptsize, figR, anchor=south west] at (3.28,0.88) {$B(y,\varepsilon)$};
\end{tikzpicture}
